\SourceLeaf{77}%
\begin{center}
{\large THEORIE DER BIQUADRATISCHEN RESTE.\par}
\bigskip ERSTE ABHANDLUNG.
\end{center}
\BookRule{28mm}

\Article{1}
Die Theorie der quadratischen Reste lässt sich auf wenige grundlegende Sätze zurückführen, die zu den schönsten Kleinodien der höheren Arithmetik gehören. Bekanntlich wurden sie zunächst durch induktive Beobachtung leicht entdeckt und dann auf vielfältige Weise so bewiesen, dass nichts mehr zu wünschen übrig bleibt.

Die Theorie der kubischen und biquadratischen Reste verlangt dagegen eine sehr viel tiefere Untersuchung. Als wir 1805 begannen, sie zu erforschen, fanden wir neben den Ergebnissen, die gleichsam an der Schwelle liegen, einige besondere Sätze, die sowohl durch ihre Einfachheit als auch durch die Schwierigkeit ihrer Beweise sehr bemerkenswert sind. Bald erkannten wir jedoch, dass die bis dahin verwendeten Prinzipien der Arithmetik keineswegs ausreichen, um die allgemeine Theorie zu begründen. Diese verlangt vielmehr notwendig, das Gebiet der höheren Arithmetik gleichsam unendlich viele Male zu erweitern. Was darunter zu verstehen ist, wird in der Fortsetzung dieser Untersuchungen völlig deutlich werden. Sobald wir dieses neue Gebiet betreten hatten, eröffnete uns die induktive Beobachtung sogleich einen Zugang zu den einfachsten Sätzen, die die gesamte Theorie umfassen. Ihre Beweise lagen aber so tief verborgen, dass sie erst nach vielen vergeblichen Versuchen schließlich ans Licht gebracht werden konnten.

Da wir nun daran gehen, diese Studien zu veröffentlichen, beginnen wir mit der Theorie der biquadratischen Reste. In dieser ersten Abhandlung\space\SourceLeaf{78}%
 werden wir diejenigen Untersuchungen darlegen, die sich schon diesseits des erweiterten Gebiets der Arithmetik abschließen ließen. Sie bahnen gleichsam den Weg dorthin und tragen zugleich einige neue Entwicklungen zur Theorie der Kreisteilung bei.

\Article{2}
Den Begriff des biquadratischen Restes haben wir in Artikel 115 der \emph{Disquisitiones Arithmeticae} eingeführt: Eine ganze Zahl $a$, positiv oder negativ, heißt biquadratischer Rest der ganzen Zahl $p$, wenn $a$ modulo $p$ einer vierten Potenz kongruent sein kann, und biquadratischer Nichtrest, wenn keine solche Kongruenz besteht. In allen folgenden Untersuchungen setzen wir, sofern wir nicht ausdrücklich etwas anderes sagen, voraus, dass der Modul $p$ eine ungerade positive Primzahl ist und dass $a$ nicht durch $p$ teilbar ist, da sich alle übrigen Fälle sehr leicht auf diesen zurückführen lassen.

\Article{3}
Offenbar ist jeder biquadratische Rest von $p$ auch quadratischer Rest derselben Zahl, und somit jeder quadratische Nichtrest auch biquadratischer Nichtrest. Dieser Satz lässt sich auch umkehren, wenn $p$ eine Primzahl der Form $4n+3$ ist. Ist nämlich in diesem Fall $a$ ein quadratischer Rest von $p$, so setzen wir $a\equiv bb\; (\mathrm{mod.}\,p)$. Dabei ist $b$ entweder quadratischer Rest von $p$ oder Nichtrest. Im ersten Fall setzen wir $b\equiv cc$, woraus $a\equiv c^4$ folgt; also ist $a$ ein biquadratischer Rest von $p$. Im zweiten Fall ist $-b$ ein quadratischer Rest von $p$ (denn $-1$ ist Nichtrest jeder Primzahl der Form $4n+3$). Setzt man $-b\equiv cc$, so folgt wie zuvor $a\equiv c^4$, und $a$ ist ein biquadratischer Rest von $p$. Zugleich sieht man leicht, dass die Kongruenz $x^4\equiv a\; (\mathrm{mod.}\,p)$ in diesem Fall außer $x\equiv c$ und $x\equiv-c$ keine weiteren Lösungen hat. Da diese unmittelbar einleuchtenden Sätze die gesamte Theorie der biquadratischen Reste für Primzahlmoduln der Form $4n+3$ erschöpfen, schließen wir solche Moduln ganz aus unserer Untersuchung aus und beschränken diese auf Primzahlmoduln der Form $4n+1$.

\Article{4}
Ist also $p$ eine Primzahl der Form $4n+1$, so lässt sich der Satz des vorigen Artikels nicht umkehren: Es kann quadratische Reste geben, die nicht zugleich biquadratische Reste sind. Dies geschieht, wenn ein quadratischer Rest dem Quadrat eines quadratischen Nichtrestes kongruent ist. Setzt man nämlich $a\equiv bb$, wobei $b$ ein quadratischer \SourceLeaf{79}%
Nichtrest von $p$ ist, und könnte die Kongruenz $x^4\equiv a$ durch einen Wert $x\equiv c$ erfüllt werden, so wäre $c^4\equiv bb$, also das Produkt $(cc-b)(cc+b)$ durch $p$ teilbar. Dann müsste $p$ entweder den Faktor $cc-b$ oder den anderen Faktor $cc+b$ teilen; das heißt, entweder $+b$ oder $-b$ wäre quadratischer Rest von $p$ und folglich wären es beide (da $-1$ ein quadratischer Rest ist), entgegen der Voraussetzung.

Man könnte also alle ganzen Zahlen, die nicht durch $p$ teilbar sind, auf drei Klassen verteilen: Die erste enthielte die biquadratischen Reste, die zweite diejenigen biquadratischen Nichtreste, die zugleich quadratische Reste sind, und die dritte die quadratischen Nichtreste. Offenbar genügt es, nur die Zahlen $1,2,3\dotsfive p-1$ so einzuteilen. Die Hälfte würde zur dritten Klasse gehören, während sich die andere Hälfte auf die erste und zweite Klasse verteilen würde.

\newpage\Article{5}
Es ist jedoch besser, vier Klassen folgendermaßen festzulegen.

Es sei $A$ die Sammlung aller biquadratischen Reste von $p$ zwischen $1$ und $p-1$, einschließlich der Grenzen, und $e$ ein beliebig gewählter quadratischer Nichtrest von $p$. Ferner sei $B$ die Sammlung der kleinsten positiven Reste der Produkte $eA$ modulo $p$, und $C,D$ seien entsprechend die Sammlungen der kleinsten positiven Reste der Produkte $eeA$, $e^3A$ modulo $p$. Dann sind offenbar die einzelnen Zahlen in $B$ voneinander verschieden, ebenso diejenigen in $C$ und diejenigen in $D$; die Null kann unter keiner dieser Zahlen vorkommen. Ferner sind alle Zahlen in $A$ und $C$ quadratische Reste von $p$, alle Zahlen in $B$ und $D$ hingegen quadratische Nichtreste. Daher kann keine der Sammlungen $A,C$ eine Zahl mit $B$ oder $D$ gemeinsam haben. Aber auch $A$ und $C$ können keine Zahl gemeinsam haben, ebenso wenig $B$ und $D$. Nehmen wir nämlich an,

I. eine Zahl aus $A$, etwa $a$, komme auch in $C$ vor und sei dort aus dem ihr kongruenten Produkt $eea'$ entstanden, wobei $a'$ eine Zahl aus $A$ ist. Man setze $a\equiv\alpha^4$, $a'\equiv\alpha'^4$ und wähle eine ganze Zahl $\Theta$ so, dass $\Theta\alpha'\equiv1$ gilt. Dann ist $ee\alpha'^4\equiv\alpha^4$, und durch Multiplikation mit $\Theta^4$ folgt

\[
ee\equiv\alpha^4\Theta^4
\]

Also wäre $ee$ ein biquadratischer Rest und somit $e$ ein quadratischer Rest, entgegen der Voraussetzung.

\SourceLeaf{80}%
II. Nimmt man ebenso an, eine Zahl gehöre zu beiden Sammlungen $B,D$ und sei aus den Produkten $ea$, $e^3a'$ entstanden, wobei $a,a'$ Zahlen aus $A$ sind, so würde aus der Kongruenz $ea\equiv e^3a'$ die Kongruenz $a\equiv eea'$ folgen. Man hätte also eine Zahl, die aus dem Produkt $eea'$ entsteht und sowohl zu $C$ als auch zu $A$ gehört, was wir gerade als unmöglich erwiesen haben.

Ferner lässt sich leicht beweisen, dass \emph{alle} quadratischen Reste von $p$ zwischen $1$ und $p-1$, einschließlich der Grenzen, notwendig entweder in $A$ oder in $C$ vorkommen und alle quadratischen Nichtreste von $p$ zwischen diesen Grenzen notwendig entweder in $B$ oder in $D$. Denn

I. Jeder solche quadratische Rest, der zugleich biquadratischer Rest ist, liegt nach Voraussetzung in $A$.

II. Einen quadratischen Rest $h$ (kleiner als $p$), der zugleich biquadratischer Nichtrest ist, schreibe man als $\equiv gg$, wobei $g$ ein quadratischer Nichtrest ist. Man wähle eine ganze Zahl $\gamma$ so, dass $e\gamma\equiv g$ gilt. Dann ist $\gamma$ ein quadratischer Rest von $p$, den wir als $\equiv kk$ schreiben. Daraus folgt

\[
h\equiv gg\equiv ee\gamma\gamma\equiv eek^4
\]

Da der kleinste Rest von $k^4$ in $A$ liegt, muss die Zahl $h$, die durch Multiplikation dieses Restes mit $ee$ entsteht, in $C$ liegen.

III. Es bezeichne $h$ einen quadratischen Nichtrest von $p$ zwischen $1$ und $p-1$. Man ermittle zwischen denselben Grenzen eine ganze Zahl $g$ so, dass $eg\equiv h$ gilt. Dann ist $g$ ein quadratischer Rest und liegt daher entweder in $A$ oder in $C$. Im ersten Fall findet sich $h$ offenbar unter den Zahlen $B$, im zweiten unter den Zahlen $D$.

Aus alledem folgt, dass sich sämtliche Zahlen $1,2,3\dotsfive p-1$ so auf die vier Listen $A,B,C,D$ verteilen, dass jede in einer davon vorkommt; jede Liste muss daher $\tfrac14(p-1)$ Zahlen enthalten. Bei dieser Einteilung stehen die Zahlen, die zu $A$ und $C$ gehören, durch das Wesen dieser Klassen fest. Die Unterscheidung zwischen $B$ und $D$ ist dagegen insofern beliebig, als sie von der Wahl der Zahl $e$ abhängt, die selbst immer zu $B$ gehört. Wählt man an ihrer Stelle eine Zahl aus $D$, so werden die Klassen $B,D$ vertauscht.

\Article{6}
Da $-1$ ein quadratischer Rest von $p$ ist, setzen wir $-1\equiv ff\; (\mathrm{mod.}\,p)$. Die vier Wurzeln der Kongruenz $x^4\equiv1$ sind dann $1,f,-1,-f$. Ist also\space%
\SourceLeaf{81}%
a ein biquadratischer Rest von $p$, etwa $\equiv\alpha^4$, so sind die vier Wurzeln der Kongruenz $x^4\equiv a$ die Zahlen $\alpha,f\alpha,-\alpha,-f\alpha$, die offenbar paarweise inkongruent sind. Sammelt man also die kleinsten positiven Reste der vierten Potenzen $1,16,81,256\dotsfour(p-1)^4$, so sind immer je vier gleich. Man erhält somit $\tfrac14(p-1)$ verschiedene biquadratische Reste, die die Sammlung $A$ bilden. Sammelt man die kleinsten Reste der vierten Potenzen nur bis $(\tfrac12p-\tfrac12)^4$, so kommt jeder zweimal vor.

\Article{7}
Das Produkt zweier biquadratischer Reste ist offenbar ein biquadratischer Rest: Multipliziert man zwei Zahlen der Klasse $A$, so gehört der kleinste positive Rest ihres Produkts wieder zu dieser Klasse. Ebenso liegen die kleinsten Reste der Produkte einer Zahl aus $B$ mit einer Zahl aus $D$ oder einer Zahl aus $C$ mit einer Zahl aus $C$ in $A$.

In $B$ liegen die Reste der Produkte $A.B$ und $C.D$; in $C$ diejenigen der Produkte $A.C$, $B.B$ und $D.D$; und schließlich in $D$ diejenigen der Produkte $A.D$ und $B.C$.

Die Beweise sind so unmittelbar, dass es genügt, einen anzudeuten. Es seien etwa $c$ und $d$ Zahlen aus $C$ und $D$, und es sei $c\equiv eea$, $d\equiv e^3a'$, wobei $a,a'$ Zahlen aus $A$ bezeichnen. Dann ist $e^4aa'$ ein biquadratischer Rest, sein kleinster Rest gehört also zu $A$. Für das Produkt gilt $cd$\relax$\equiv e.e^4aa'$; sein kleinster Rest liegt daher in $B$.

Nun lässt sich auch leicht entscheiden, zu welcher Klasse ein Produkt mehrerer Faktoren gehört. Gibt man den Klassen $A,B,C,D$ entsprechend die Charaktere $0,1,2,3$, so ist der Charakter des Produkts entweder die Summe der Charaktere der einzelnen Faktoren oder deren kleinster Rest modulo $4$.

\Article{8}
Es schien lohnend, diese elementaren Sätze ohne die Theorie der Potenzreste zu entwickeln. Zieht man diese heran, so lässt sich alles noch sehr viel leichter beweisen.

\Needspace{16\baselineskip}Es sei $g$ eine primitive Wurzel für den Modul $p$, also eine Zahl, bei der in der Folge der Potenzen $g,gg,g^3\dotsfour$ keine vor $g^{p-1}$ modulo $p$ der Einheit kongruent ist. Dann stimmen die kleinsten positiven Reste der Zahlen $1,g,gg,g^3\dotsfour g^{p-2}$, von ihrer Reihenfolge abgesehen, mit $1,2,3\dotsfive p-1$ überein und verteilen sich folgendermaßen auf vier Klassen:

\SourceLeaf{82}%
\[
\begin{array}{c|l}
\text{zu}&\text{kleinste Reste der Zahlen}\\\hline
A&1,\ g^4,g^8,g^{12}\dotsfive g^{p-5}\\
B&g,\ g^5,g^9,g^{13}\dotsfive g^{p-4}\\
C&gg,g^6,g^{10},g^{14}\dotsfive g^{p-3}\\
D&g^3,g^7,g^{11},g^{15}\dotsfive g^{p-2}
\end{array}
\]

Daraus ergeben sich alle vorhergehenden Sätze unmittelbar.

So wie hier die Zahlen $1,2,3\dotsfour p-1$ auf vier Klassen verteilt wurden, deren Sammlungen wir mit $A,B,C,D$ bezeichnen, kann man \emph{jede} ganze Zahl, die nicht durch $p$ teilbar ist, nach ihrem kleinsten Rest modulo $p$ einer dieser Klassen zuordnen.

\Article{9}
Wir bezeichnen mit $f$ den kleinsten Rest der Potenz $g^{\frac14(p-1)}$ modulo $p$. Da $ff\equiv g^{\frac12(p-1)}\equiv-1$ gilt (\emph{Disquis. Arithm.}, Artikel 62), hat das Zeichen $f$ hier offenbar dieselbe Bedeutung wie in Artikel 6. Die Potenz $g^{\frac14\lambda(p-1)}$, wobei $\lambda$ eine positive ganze Zahl bezeichnet, ist also modulo $p$ der Zahl $1,f,-1,-f$ kongruent, je nachdem $\lambda$ entsprechend die Form $4m,4m+1,4m+2,4m+3$ hat oder der kleinste Rest von $g^\lambda$ entsprechend in $A,B,C,D$ liegt. Dies gibt ein sehr einfaches Kriterium dafür, welcher Klasse eine gegebene Zahl $h$, die nicht durch $p$ teilbar ist, zuzuordnen ist: $h$ gehört zu $A,B,C$ oder $D$, je nachdem die Potenz $h^{\frac14(p-1)}$ modulo $p$ der Zahl $1,f,-1$ oder $-f$ kongruent ist.

Als Folgerung ergibt sich, dass $-1$ stets zur Klasse $A$ gehört, wenn $p$ die Form $8n+1$ hat, und zur Klasse $C$, wenn $p$ die Form $8n+5$ hat. Ein von der Theorie der Potenzreste unabhängiger Beweis dieses Satzes lässt sich leicht aus dem gewinnen, was wir in Artikel 115, III der \emph{Disquisitiones Arithmeticae} gezeigt haben.

\Article{10}
Da \emph{alle} primitiven Wurzeln für den Modul $p$ aus den Resten der Potenzen $g^\lambda$ hervorgehen, wenn man für $\lambda$ alle zu $p-1$ teilerfremden Zahlen nimmt, sieht man leicht, dass diese Wurzeln gleichmäßig auf $B$ und $D$ verteilt sind, wobei die Basis $g$ stets in $B$ liegt. Wählt man anstelle von $g$ eine andere primitive Wurzel aus $B$ als Basis, so bleibt die Einteilung dieselbe; wählt man dagegen eine primitive Wurzel aus $D$ als Basis, so werden die Klassen $B$ und $D$ vertauscht.

\SourceLeaf{83}%
Beruht die Einteilung auf dem Kriterium des vorigen Artikels, so hängt die Unterscheidung zwischen $B$ und $D$ davon ab, welche Wurzel der Kongruenz $xx\equiv-1\; (\mathrm{mod.}\,p)$ wir als charakteristische Zahl $f$ wählen.

\Article{11}
Damit sich die feineren Untersuchungen, die wir nun beginnen, leichter durch Beispiele veranschaulichen lassen, geben wir hier die Konstruktion der Klassen für alle Moduln unter $100$ an. Für jeden haben wir die kleinste primitive Wurzel gewählt.

\[
\begin{gathered}
p=5\\
g=2,\ f=2\\[3pt]
\begin{array}{c|l}
A&1\\
B&2\\
C&4\\
D&3
\end{array}
\end{gathered}
\]

\[
\begin{gathered}
p=13\\
g=2,\ f=8\\[3pt]
\begin{array}{c|l}
A&1,3,9\\
B&2,5,6\\
C&4,10,12\\
D&7,8,11
\end{array}
\end{gathered}
\]

\[
\begin{gathered}
p=17\\
g=3,\ f=13\\[3pt]
\begin{array}{c|l}
A&1,4,13,16\\
B&3,5,12,14\\
C&2,8,9,15\\
D&6,7,10,11
\end{array}
\end{gathered}
\]

\[
\begin{gathered}
p=29\\
g=2,\ f=12\\[3pt]
\begin{array}{c|l}
A&1,7,16,20,23,24,25\\
B&2,3,11,14,17,19,21\\
C&4,5,6,9,13,22,28\\
D&8,10,12,15,18,26,27
\end{array}
\end{gathered}
\]

\SourceLeaf{84}%
\[
\begin{gathered}
p=37\\
g=2,\ f=31\\[3pt]
\begin{array}{c|l}
A&1,7,9,10,12,16,26,33,34\\
B&2,14,15,18,20,24,29,31,32\\
C&3,4,11,21,25,27,28,30,36\\
D&5,6,8,13,17,19,22,23,35
\end{array}
\end{gathered}
\]

\[
\begin{gathered}
p=41\\
g=6,\ f=32\\[3pt]
\begin{array}{c|l}
A&1,4,10,16,18,23,25,31,37,40\\
B&6,14,15,17,19,22,24,26,27,35\\
C&2,5,8,9,20,21,32,33,36,39\\
D&3,7,11,12,13,28,29,30,34,38
\end{array}
\end{gathered}
\]

\[
\begin{gathered}
p=53\\
g=2,\ f=30\\[3pt]
\begin{array}{c|l}
A&1,10,13,15,16,24,28,36,42,44,46,47,49\\
B&2,3,19,20,26,30,31,32,35,39,41,45,48\\
C&4,6,7,9,11,17,25,29,37,38,40,43,52\\
D&5,8,12,14,18,21,22,23,27,33,34,50,51
\end{array}
\end{gathered}
\]

\[
\begin{gathered}
p=61\\
g=2,\ f=11\\[3pt]
\begin{array}{c|l}
A&1,9,12,13,15,16,20,22,25,34,42,47,56,57,58\\
B&2,7,18,23,24,26,30,32,33,40,44,50,51,53,55\\
C&3,4,5,14,19,27,36,39,41,45,46,48,49,52,60\\
D&6,8,10,11,17,21,28,29,31,35,37,38,43,54,59
\end{array}
\end{gathered}
\]

\[
\begin{gathered}
p=73\\
g=5,\ f=27\\[3pt]
\begin{array}{c|l}
A&1,2,4,8,9,16,18,32,36,37,41,55,57,64,65,69,71,72\\
B&5,7,10,14,17,20,28,33,34,39,40,45,53,56,59,63,66,68\\
C&3,6,12,19,23,24,25,27,35,38,46,48,49,50,54,61,67,70\\
D&11,13,15,21,22,26,29,30,31,42,43,44,47,51,52,58,60,62
\end{array}
\end{gathered}
\]

\SourceLeaf{85}%
\[
\begin{gathered}
p=89\\
g=3,\ f=34\\[3pt]
\begin{array}{c|l}
A&1,2,4,8,11,16,22,25,32,39,44\mathpunct{.}\,45,\\
&50,57,64,67,73,78,81,85,87,88\\
B&3,6,7,12,14,23,24,28,33,41,43,46,\\
&48,56,61,65,66,75,77,82,83,86\\
C&5,9,10,17,18,20,21,34,36,40,42,47,\\
&49,53,55,68,69,71,72,79,80,84\\
D&13,15\mathpunct{.}\,19,26,27,29\mathpunct{.}\,30,31,35,37,38\mathpunct{.}\,51,\\
&52,54,58,59,60,62,63,70,74,76
\end{array}
\end{gathered}
\]

\[
\begin{gathered}
p=97\\
g=5,\ f=22\\[3pt]
\begin{array}{c|l}
A&1,4,6,9,16,22,24,33,35,36,43,47,\\
&50,54,61,62,64,73,75,81,88,91,93,96\\
B&5,13,14,17,19,20,21,23,29,30,41,45,\\
&52,56,67,68,74,76,77,78,80,83,84,92\\
C&2,3,8,11,12,18,25,27,31,32,44,48,\\
&49,53,65\mathpunct{.}\,66,70,72,79,85,86,89,94,95\\
D&7,10,15,26,28,34,37,38,39\mathpunct{.}\,40,42,46,\\
&51,55,57,58,59,60,63,69,71,82,87,90
\end{array}
\end{gathered}
\]

\Article{12}
Da $2$ quadratischer Rest aller Primzahlen der Form $8n+1$, hingegen Nichtrest aller Primzahlen der Form $8n+5$ ist, liegt für Primzahlmoduln der ersten Form $2$ in der Klasse $A$ oder $C$, für Moduln der zweiten Form dagegen in $B$ oder $D$. Der Unterschied zwischen $B$ und $D$ ist nicht wesentlich, da er allein von der Wahl der Zahl $f$ abhängt; daher lassen wir die Moduln der Form $8n+5$ vorerst beiseite. Durch eine \emph{induktive Untersuchung} der Moduln der Form $8n+1$ finden wir, dass $2$ zu $A$ gehört für $p=73,89,113,233,257,281,337,353$ usw.; dagegen gehört $2$ zu $C$ für $p=17,41,97,137,193,241,313,401,409,433,449,457$ usw.

Da ferner für einen Primzahlmodul der Form $8n+1$ die Zahl $-1$ ein biquadratischer Rest ist, gehört $-2$ offenbar stets zur selben Klasse wie $+2$.
